\section{Methodisches Vorgehen}

\subsection{Orientierung an Design Science Research}
DSRM Entwurf:
1. Problemidentifikation und Motivation: 
    Risiken einer unreflektierten KI-Nutzung beim Kompetenzaufbau aufzeigen und Handlungsbedarf für Unternehmen begründen
2. Zieldefinition:
    Anforderungen an den Leitfaden aus Forschungsliteratur und ergänzender Anfängerbefragung ableiten
3. Design und Entwicklung:
    Wissenschaftlich begründeten Leitfaden mit KI-Nutzungsregeln und Kontrollmöglichkeiten entwickeln
4. Demonstration 
    	Leitfaden durch Programmieranfänger bei einer beispielhaften Programmieraufgabe anwenden lassen
5. Evaluation
    	Anwendbarkeit, Verständlichkeit und subjektive Lernunterstützung durch Anfänger sowie fachliche Eignung durch Experten (programmierdozent) bewerten lassen
6. Dokumentation/Kommunikation
    Vorgehen und Limitationen dokumentieren


Die Konzeption des Leitfadens orientiert sich an der Design Science Research Methodology (DSRM) nach Peffers et al. (2007). Diese beschreibt ein strukturiertes Vorgehen zur Entwicklung und Evaluation von Artefakten, die der Lösung identifizierter Probleme dienen. Das Vorgehensmodell umfasst sechs Aktivitäten: Problemidentifikation (Bei mir in Kapitel 5.2 und 5.3) und Motivation, sowie Zieldefinition (beides in Einleitung), Design und Entwicklung (Kapitel 6), Demonstration (Kapitel 7), Evaluation sowie Kommunikation (Kapitel 7) (vgl. Peffers et al., 2007, S. 12–14).

Im ersten Schritt der DSRM erfolgt die Identifikation und Begründung der zugrunde liegenden Problemstellung. Hierzu werden zunächst mögliche Auswirkungen des Einsatzes generativer KI auf den Lernprozess von Programmieranfängern anhand wissenschaftlicher Literatur untersucht (Kapitel 5.2). Anschließend werden die potenziellen Risiken unzureichender Programmierkompetenzen im Unternehmenskontext betrachtet (Kapitel 5.3).
Ergänzend dient eine empirische Onlinebefragung dazu, das KI-Nutzungsverhalten von Programmieranfängern sowie deren subjektive Wahrnehmung des Einflusses auf den eigenen Lernprozess zu erfassen. Dadurch wird die literaturbasierte Problemanalyse um die Perspektive der Lernenden erweitert. Die methodische Konzeption der Befragung wird in Kapitel 2.3 beschrieben.

Auf Grundlage der wissenschaftlichen Erkenntnisse zu wichtigen Lernimpulsen bei Programmieranfängern werden anschließend Anforderungen an einen Leitfaden zur lernförderlichen KI-Nutzung abgeleitet. Diese bilden die Grundlage für die Konzeption und Entwicklung des Leitfadens als methodisches Artefakt.

Die anschließende Demonstration erfolgt durch eine exemplarische Anwendung des Leitfadens bei der Bearbeitung einer Programmieraufgabe durch Programmieranfänger. Die anschließende Evaluation soll dessen Verständlichkeit, praktische Anwendbarkeit und subjektiv wahrgenommene Lernunterstützung untersuchen. Ergänzend ist eine Expertenbewertung zur Beurteilung der fachlichen und didaktischen Eignung vorgesehen.

Aufgrund des begrenzten Bearbeitungszeitraums erfolgt die Evaluation lediglich in eingeschränktem Umfang. Insbesondere können langfristige Auswirkungen auf den Kompetenzaufbau sowie eine kausale Wirksamkeit gegenüber anderen Formen der KI-Nutzung nicht nachgewiesen werden. Die Untersuchung konzentriert sich daher auf eine erste Bewertung der Eignung und praktischen Anwendbarkeit des entwickelten Leitfadens. Die Ergebnisse sollen zur Identifikation möglicher Verbesserungsbedarfe herangezogen werden.

\subsection{Vorgehen bei der Literaturrecherche}
Die Literaturrecherche orientiert sich am Vorgehensmodell nach vom Brocke et al. (2009), welches die Definition des Untersuchungsumfangs, die Konzeptualisierung des Themengebiets, die Literatursuche sowie die Analyse und Synthese der identifizierten Veröffentlichungen umfasst.

Zunächst werden die für die Forschungsfragen relevanten Themenbereiche und zentralen Suchbegriffe festgelegt. Anschließend erfolgt die Recherche in wissenschaftlichen Datenbanken anhand geeigneter Suchbegriffe und Suchkombinationen. Ergänzend sollen durch Rückwärts- und Vorwärtssuchen weitere relevante Veröffentlichungen identifiziert werden.

Die Auswahl der Literatur erfolgt anhand ihrer thematischen Relevanz für die Forschungsfragen und ihrer wissenschaftlichen Qualität. Die gewonnenen Erkenntnisse werden thematisch geordnet, miteinander verglichen und zur Ableitung der Anforderungen an den zu entwickelnden Leitfaden herangezogen.

Zur Gewährleistung der Nachvollziehbarkeit werden die verwendeten Suchbegriffe, Datenbanken und Auswahlentscheidungen dokumentiert.

\subsection{Vorgehen bei der Fragebogenentwicklung}

Zur Ergänzung der literaturbasierten Untersuchung ist eine Onlinebefragung unterstützend für die Beantwortung der im Zuge der zweiten Forschungsfrage in Kapitel 5.2 gestellte Frage nach den Auswirkungen von KI-Nutzung auf den Lernprozess vorgesehen. 
Ziel ist es, das Nutzungsverhalten generativer KI bei Programmieranfängern sowie deren subjektiv wahrgenommenen Einfluss auf den eigenen Lernprozess zu untersuchen. Dabei geht es hauptsächlich um die Auswirkungen auf das unmittelbare Codeverständnis, die Fähigkeit zur eigenen Problemlösung und den langfristigen Kompetenzaufbau in der Softwareentwicklung.

Die Untersuchung ist als deskriptiv-explorative Querschnittsbefragung konzipiert. Während die deskriptive Komponente der Beschreibung des KI-Nutzungsverhaltens und der wahrgenommenen Auswirkungen dient, sollen im explorativen Teil mögliche Zusammenhänge zwischen Programmiererfahrung, Nutzungsweisen und subjektiv wahrgenommenen Lerneinflüssen untersucht werden. Die Datenerhebung erfolgt einmalig mittels eines standardisierten Online-Fragebogens.

Die methodische Konzeption orientiert sich an den Richtlinien zur Durchführung von Umfragen im Software Engineering. Molléri et al. (2016) identifizieren und vergleichen hierzu bestehende methodische Vorgehensmodelle. Als Grundlage für die geplante Erhebung dienen insbesondere die von Kasunic (2005) und Linåker et al. (2015) beschriebenen Vorgehensweisen.

Die Entwicklung des Fragebogens soll sich am Goal-Question-Metric-Ansatz orientieren, der bei Linåker et al. (2015, S. 10) beschrieben wird. Dabei werden aus den Untersuchungszielen zunächst übergeordnete Fragestellungen und anschließend messbare Merkmale sowie konkrete Fragebogenitems abgeleitet.

Der Fragebogen soll vier Themenbereiche umfassen: Teilnahmevoraussetzungen, Ausbildung und Programmiererfahrung, KI-Nutzungsverhalten sowie subjektiv wahrgenommene Auswirkungen auf den Lernprozess.

Die Rekrutierung der Teilnehmenden ist über universitäre E-Mail-Verteiler sowie betriebliche Kontakte zu Studierenden und Auszubildenden vorgesehen. Aufgrund der freiwilligen Teilnahme und der eingeschränkten Verfügbarkeit der potenziell Teilnehmenden handelt es sich um eine nicht-probabilistische Stichprobe. Die damit verbundenen Einschränkungen der Übertragbarkeit auf die gesamte Zielgruppe wird bei der Interpretation der Ergebnisse berücksichtigt.

Vor der eigentlichen Datenerhebung ist eine Überprüfung des Fragebogens durch einen Pretest mit Personen aus der Zielgruppe vorgesehen. Dabei sollen insbesondere Verständlichkeit, Eindeutigkeit und technische Funktion überprüft werden, wenn möglich über einige kognitive Interviews

\subsection{Vorgehen bei der Evaluation des Leitfadens}
